\title{LongRoPE: Extending LLM Context Window Beyond 2 Million Tokens}

\begin{abstract}
Expanding the context window is a highly sought-after attribute in large language models (LLMs). Nevertheless, the reach of present-day extended context windows remains constrained to approximately 128k tokens, primarily due to the substantial expenses associated with fine-tuning, the limited availability of lengthy textual data, and the destabilizing effects caused by introducing novel token positions.
\end{abstract}

This work presents LongRoPE, a method that, for the first time, extends the context window of pre-trained LLMs to an extraordinary \textbf{2048k} tokens, requiring no more than 1k fine-tuning steps at training lengths up to 256k, all while preserving performance at the initial short context size. This advancement relies on three primary innovations: (i) we detect and utilize two distinct forms of non-uniformities in positional interpolation using an efficient search procedure, resulting in improved fine-tuning initialization and enabling up to an 8$\times$ increase in context length without fine-tuning; (ii) we design a staged extension protocol in which the LLM is first fine-tuned on 256k token sequences, followed by an additional round of positional interpolation on this extended model, thereby unlocking a 2048k context window; (iii) we recalibrate LongRoPE on a length of 8k tokens in order to restore the model’s original short context window capabilities. Comprehensive experiments using LLaMA2 and Mistral on a range of benchmarks validate the superiority of our approach. The models adapted with LongRoPE keep the original model design unchanged except for minor adjustments to the positional embedding, and can take advantage of most established optimizations. The implementation will be released at \texttt{https://github.com/microsoft/LongRoPE}

\section{Introduction}

Although Large Language Models (LLMs) have demonstrated impressive performance across a range of tasks~\cite{gpt4,llama2}, their context windows are inherently restricted; for instance, LLaMA2 can only process up to 4096 tokens at once~\cite{llama2}. When inputs exceed this window, LLM effectiveness deteriorates, as the models lack exposure to such extended positional sequences during training. This limitation becomes problematic in critical use cases, such as in-context learning involving many examples~\cite{Huang2023FewerIM}, and the deployment of LLM agents~\cite{park2023generative,madaan2023selfrefine}.

Recent studies demonstrate that the context window in a pre-trained LLM can be pushed to roughly 128k tokens through fine-tuning with longer input sequences~\cite{longlora,pi,yarn,zhang2024soaring,liu2023scaling}. However, three primary challenges hinder the possibility of increasing this limit further. Firstly, newly added position indices that lack prior training can produce numerous anomalous values. This results in out-of-distribution complications and makes the fine-tuning process difficult to reach convergence~\cite{pi}. The situation becomes especially problematic during an expansion from 4k to more than 1000k tokens, as over 90\% of the position indices are novel. Secondly, suitable textual data matching these vast sequence lengths is often necessary for fine-tuning, but existing datasets rarely include samples that exceed 1000k tokens. Additionally, processing such lengthy sequences is highly resource-intensive, demanding excessive training time and significant GPU capacity. Thirdly, for exceptionally large context windows, the attention mechanism spreads its focus across a multitude of token positions, diluting its effectiveness and diminishing performance when handling sequences of shorter length~\cite{pi}.

To address the initial challenge, a commonly used technique is interpolating RoPE positional embeddings~\cite{rope,pi}, which involves scaling down novel positional indices to fit within the range established during pretraining, as depicted in Fig.\ref{fig:overview}. The Position Interpolation (PI) method~\cite{pi} implements a linear adjustment to RoPE's rotary angles, proportional to the extension ratio. NTK~\cite{ntk,dynamicntk} proposes a method with non-uniform interpolation and extrapolation along RoPE dimensions. YaRN~\cite{yarn} takes a frequency-based approach, segmenting RoPE dimensions into three groups and applying extrapolation, NTK, or linear interpolation, accordingly. Despite these innovations, positional embeddings in Transformers demonstrate intricate, non-uniform information entropy. Existing techniques fail to sufficiently exploit these nuanced variations, resulting in a loss of information that ultimately constrains the maximum achievable context window size.

Section~\ref{sec:analysis} provides two principal empirical observations: \textbf{(1)} Successful positional interpolation must account for two types of non-uniformity—differences across RoPE dimensions and variations in token positions. Less interpolation suffices for smaller RoPE dimensions and for tokens at the sequence’s beginning; however, the most suitable approach will differ depending on the intended extended context length.
\textbf{(2)} Incorporating these sources of non-uniformity into positional interpolation strategies enables preservation of critical information from the original RoPE encoding, especially in specific dimensions and token positions. This approach effectively reduces the loss inherent to positional interpolation, thereby offering superior initialization for downstream fine-tuning tasks. In addition, it supports up to an 8$\times$ context extension even without further fine-tuning.

Building upon these results, we present \textbf{LongRoPE}, a powerful approach for enlarging the LLM context window beyond 2 \emph{million} tokens. The design of LongRoPE rests on three principal contributions. To begin with, {LongRoPE} maximally utilizes multidimensional variations within positional interpolation, pinpointing optimal rescale factors for RoPE's rotational angles on a per-dimension basis and depending on the locations of tokens. Because the task of finding suitable rescale factors grows exponentially in complexity relative to the intended window expansion, {LongRoPE} adopts an evolutionary search procedure equipped with two optimizations aimed at enhancing the search process. An illustrative case of the optimized rescaled RoPE is provided in Fig.~\ref{fig:overview}.

Subsequently, {LongRoPE} employs a resource-efficient, stepwise extension method to reach a 2048k context window, circumventing the need for direct fine-tuning on ultra-long text sequences, which are limited in availability. This process initiates by identifying an optimal 256k context length using the original pre-trained LLM, followed by fine-tuning the model specifically for this length. Leveraging the capacity of our non-uniform positional interpolation to extend context up to 8$\times$ in scenarios lacking fine-tuning, we then execute a secondary search to determine updated RoPE rescale factors using the already fine-tuned, extended LLM. This methodology allows us to realize a 2048k context window for both LLaMA2 and Mistral~\cite{mistral}.

In conclusion, to address potential decreases in performance within the initial (shorter) context window, {LongRoPE} persistently tunes the RoPE rescaling parameters on the expanded LLM. Following an analogous procedure to extending the context window from 256k to 2048k, we apply a downscaling process to the 256k fine-tuned LLM, adjusting it to accommodate 4k and 8k context windows via our optimization algorithm, thereby minimizing excessive positional interpolation. At inference time, for input sequences shorter than 8k, RoPE is revised using the rescale factors determined by the search procedure.

A wide range of experiments conducted on multiple LLMs and diverse long-context tasks validate the efficacy of our approach. Our findings indicate that LongRoPE consistently preserves low perplexity levels across evaluation lengths spanning from 4k up to 2048k, attains more than 90\% accuracy in passkey retrieval, and matches standard benchmark performance established for the 4096 context window. The LongRoPE technique is adaptable to any LLM utilizing RoPE embedding. We plan to make both our code and LongRoPE-2048k models publicly available.

\section{Non-uniformity in Positional Interpolation}
\label{sec:analysis}

\subsection{Preliminary}

Transformer models require explicit positional information, often in the form of position embedding, to represent the order of input tokens. Our work focuses on the RoPE~\cite{rope} position embedding, which is widely used in recent  LLMs. For a token at position index $n$, its corresponding RoPE encoding can be simplified as follows:

\begin{equation}
		\fontsize{7.8}{7.8} \selectfont
	\label{eq:rope}
	\left[ cos(n\theta_0), sin(n\theta_0), cos(n\theta_1), \cdot\cdot\cdot, cos(n\theta_{d/2-1}), sin(n\theta_{d/2-1}) \right]   
\end{equation}

Here, $d$ denotes the dimension of the embeddings, $n\theta_i$ specifies the rotary angle associated with the token at position $n$, and $\theta_i=\theta^{-2i/d}$ encodes the frequencies for rotation. Within RoPE, the typical setting for the base $\theta$ is 10000.

\textbf{\emph{Context window extension ratio $s$} and positional interpolation.} The context extension ratio $s$ is characterized as the quotient of the new context length $L'$ divided by the original context length $L$, such that $s=\frac{L'}{L}$.

To extend the context window from $L$ to $L'$, current positional interpolation methods suggest downscaling rotation frequencies $\theta_i$ based on extension ratio $s$. Let $\beta=\theta^{2/d}$, and $\lambda$ denote the actual rescale factor related to $s$, we   unify these positional interpolation methods as follows:

\begin{equation}	
	\fontsize{7.5}{7.5} \selectfont
	\label{eq:unified_rope}
	\begin{aligned}
		\left[cos\left(\frac{n}{\lambda(\beta)^0}\right), sin \left(\frac{n}{\lambda(\beta)^0}\right), cos\left(\frac{n}{\lambda(\beta)^1}\right), \cdot\cdot\cdot,  sin \left(\frac{n}{\lambda(\beta)^{d/2-1}}\right) \right]   
	\end{aligned}
\end{equation}

\textbf{Linear positional interpolation (PI)}. PI~\cite{pi} employs linear interpolation of positional indices as long as they fall within the pre-trained sequence length. Given a desired extension ratio $s$, it uniformly scales down the rotation angles of all positions in every RoPE dimension by a factor of $\lambda = s$. This uniform compression, however, compresses positional representations, which diminishes the model’s capacity to separate tokens that are near each other. As a result, PI typically experiences a decline in effectiveness when applied at large extension ratios.

\textbf{NTK-based interpolation and extrapolation}. The works of~\cite{ntk,dynamicntk} examine RoPE from the angle of information encoding and employ Neural Tangent Kernel (NTK) theory~\cite{jacot2018neural,tancik2020fourier} in their analysis. To address the issue of densely packed positional representations in PI, their approach redistributes the burden of interpolation over the various RoPE dimensions. This is accomplished by diminishing the scaling of dimensions with lower indices (high frequency), while amplifying scaling for higher-indexed (low frequency) dimensions, thereby facilitating both interpolation and extrapolation of positional data, with $\lambda=s^{i}$. The enhanced dynamic NTK model~\cite{dynamicntk} modifies the extension ratio dynamically at each position, in accordance with the prevailing sequence length. In contrast to PI, which demands additional fine-tuning, NTK-aware techniques can broaden context windows even in the absence of fine-tuning; however, the maximum achievable extension ratio is typically limited to 4$\times$.

\textbf{YaRN}~\cite{yarn} offers a notable advance in the effectiveness of positional interpolation. This method categorizes RoPE dimensions into three groups, based on their associated frequencies, and assigns a tailored interpolation technique to each group. Specifically, dimensions with high frequencies are subject to extrapolation ($\lambda$=1), while those with low frequencies apply linear interpolation (PI). Dimensions with intermediate frequencies are addressed using NTK. YaRN's principal innovation stems from this systematic partitioning of RoPE dimensions; however, the current approach relies on manual, empirical tuning by researchers, which could limit optimality when adapting to novel LLMs.

\subsection{Study on Non-uniform Positional Interpolation}

Drawing from insights provided by NTK and YaRN, we observe that improvements arise primarily due to the incorporation of non-linear approaches, particularly when treating frequencies differently within various RoPE dimensions to better manage interpolation and extrapolation tasks. Yet, the prevailing non-linear strategies depend largely on heuristics crafted by researchers. This observation prompts two main inquiries: (1) Does present positional interpolation achieve optimality? (2) Might there exist additional, as-yet-unexplored forms of non-linearity?

In pursuit of addressing these questions, we employ evolution search (refer to Sec.~\ref{sec:method}) to uncover improved non-uniform positional interpolations for LLaMA2-7B. The perplexity metric serves as our optimization signal, evaluated over 5 randomly drawn passages from the PG19~\cite{pg19} validation split. Our experimental investigation yields several important insights as detailed below.

\textbf{Finding 1}: \textit{RoPE dimensions exhibit substantial non-uniformities, which are not effectively handled by current positional interpolation methods.}

We perform an optimization of $\lambda$ for each RoPE dimension as described in Eq.~\ref{eq:unified_rope}. In Table~\ref{tbl:exp1}, we present the perplexity results for LLaMA2-7B across different approaches on the PG19 and Proof-pile~\cite{proof-pile} test sets, all evaluated without any additional fine-tuning. The solution identified by our search yields marked improvements, indicating that both the prevailing linear (PI) and non-uniform (Dynamic-NTK, YaRN) interpolation techniques are not optimal. Furthermore, YaRN exhibits worse performance than both PI and NTK on PG19, largely because it fails to accommodate the intended context window length in the absence of LLM fine-tuning. As an illustration, with YaRN, the perplexity sharply increases after 7,000 tokens within an 8,000-token context window.

In our investigation, we observed that the rescaled factors $\lambda$ in Eq.~\ref{eq:unified_rope} are no longer uniform, in contrast to the constant scaling $s$ applied in PI, the analytic method in NTK, or the group-based adjustment of YaRN. These varying rescaling factors yield marked improvements in LLaMA2's language modeling capability—as measured by perplexity—when processing context lengths of 8k and 16k tokens, all achieved without additional fine-tuning. This enhancement stems from the modified positional encoding's superior preservation of the structure inherent to RoPE, particularly maintaining key dimensions, which lessens the model's confusion when discerning the positions of nearby tokens.

\textbf{Finding 2}: \textit{RoPE for the initial tokens in the input sequence should be extrapolated with less interpolation.} 

For the first $\hat{n}$ tokens in the input sequence, we propose that RoPE should be less reliant on interpolation. The rationale is that these tokens attract higher attention scores, establishing their significance in the attention mechanism—an effect also discussed in Streaming LLM~\cite{streamingllm} and LM-Infinite~\cite{han2023lminfinite}. To test this idea, we increase the context window to 8k and 16k using PI and NTK methodologies, systematically leaving the initial $\hat n$ tokens (where $\hat n$ ranges from 0 to 256) unchanged by interpolation. When $\hat n$=0, the setup corresponds to the standard PI and NTK approach. Table~\ref{tbl:tokenposition} provides two main insights: \textbf{(1)} skipping position interpolation for the leading tokens effectively boosts performance, and \textbf{(2)} the ideal value for $\hat n$ varies according to the desired length of context extension.

\textbf{Finding 3}: \textit{Non-uniform positional interpolation effectively extends LLM context window in both fine-tuning and non-fine-tuning settings.} 

Although our non-uniform position interpolation strategy delivers substantial improvements in context extension at 8k and 16k without additional training, scaling to even larger context windows necessitates fine-tuning. Therefore, we conduct fine-tuning of LLaMA2-7B utilizing our optimized RoPE approach for a 64k context length (refer to the Appendix for detailed configurations). Table~\ref{tbl:finetune} demonstrates that our approach yields considerably superior results compared to PI and YaRN, both prior to and following the fine-tuning of LLaMA2-7B. This advantage stems from our robust exploitation of non-uniform positional interpolation, which reduces information loss and offers improved initialization for subsequent fine-tuning.

\textbf{Summary}. Through our investigation, we identify two sources of non-uniformity: variation in RoPE dimensions and across token positions. Effectively leveraging these non-uniform aspects in positional interpolation markedly enhances the context extension abilities of LLMs.

\section{LongRoPE}

\label{sec:method}
Building on these observations, we propose LongRoPE, which initially incorporates a search algorithm designed to efficiently harness the two identified non-uniformities. Leveraging this mechanism, LongRoPE is capable of extending the LLM context window to exceed 2 million tokens.

\subsection{Problem Formulation}

These two types of non-uniformity contribute to an expansive set of possible solutions and complicate the optimization process. To mitigate these challenges, we recast the optimization problem for multidimensional non-uniform position interpolation as a search problem.

For a LLM targeting a  context window size of $L'$ and lengthy input documents $\mathbf{X}$, where each $\mathbf{x}\in\mathbf{X}$ surpasses $L'$ in token length, we denote the original rotary angle of the $i^{th}$ dimension in RoPE embedding at token position $n$ as  $\frac{n}{\beta^i}$. The optimization problem is then formulated as follows:

\begin{equation}
\begin{gathered}
		\label{eq:problem}

\underset {\mathbf{x}\in\mathbf{X}; \, |\mathbf{x}|\ge L'}{ \operatorname {arg\,min} } \,  \mathcal{L}\, \left( \text{LLM} (\text{RoPE}, \mathbf{X})\right), \text{such that} 
\\
\scalebox{0.93}{
$\underset {\begin{subarray}{l} i=0,\cdot\cdot,\frac{d}{2}-1;\\ n\in[ 0,|\mathbf{x}|);\end{subarray}}{\text{RoPE}_(n)}= {\left[\cdots,\cos\left(\mathbb{I}(\hat{\lambda}_{i}, \hat{n})\times\frac{n}{\beta^i}\right), \sin\left(\mathbb{I}(\hat{\lambda}_{i}, \hat{n})\times\frac{n}{\beta^i}\right),\cdots\right] }$}\\
	\text{with} \mathbb{I}(\hat{\lambda}_{i},\hat{n})=
\begin{cases}
	1&\text{\;  if $n<\hat{n}$}\\
	{\frac{1}{\lambda}_{i}}&\text{\; if $n\geq\hat{n}$}
\end{cases}
	\end{gathered}
\end{equation}

In this context, a series of rescale factors, $\mathbb{I}(\hat{\lambda}_{i},\hat{n})$, are introduced to address both forms of non-uniformities present. Here, $\hat{\lambda_i}$ represents the non-uniformity within the RoPE dimensions, while $\hat{n}$ specifies the variation in token positions. The rotation angle associated with the $i^{th}$ RoPE dimension is modified using $\mathbb{I}(\hat{\lambda}_{i},\hat{n})$, wherein $\hat\lambda_{i}$ serves as the corresponding rescale coefficient and $\hat{n}$ defines the positional threshold for tokens. Up to the first $\hat{n}$-1 tokens, the rescale factor $\hat\lambda_{i}$ is not incorporated, leaving the original RoPE rotary angle $\frac{n}{\beta^i}$ intact. For tokens where $n\ge\hat{n}$, the adjustment by the rescale factor is implemented.

Suppose the desired context window size is $L'$. The task is to determine the best set of rescale factors—specifically, $\mathbb{I}(\hat{\lambda}_{0},\hat{n})$, $\mathbb{I}(\hat{\lambda}_{1},\hat{n})$, ...,$\mathbb{I}(\hat{\lambda}_{i},\hat{n})$, and so forth—covering each RoPE dimension from the first up to the $d$-th. With these optimized rescale factors applied to the $\text{RoPE}$ within the target $\text{LLM}$, the model is expected to achieve the lowest possible next token prediction loss, $\mathcal{L}$ (which reflects perplexity), when processing input data $\mathbf{X}$ of length $L'$.

\subsection{Searching the Non-uniform Position Interpolation}

To address the challenge outlined in Eq.~\ref{eq:problem}, we propose a straightforward but powerful technique that identifies the best RoPE rescale factors, enabling comprehensive utilization of the diverse multidimensional variations present in position embedding.

\textbf{Search space}. We construct an extensive search space that accommodates the two forms of non-uniformity. The configuration is summarized in Table~\ref{tbl:searchspace}. In detail, our approach enables exploration of individual rescale factors for each dimension within RoPE embedding. For clarity and manageability of the search space, we focus on optimizing $\lambda_i$ and $\hat{n}$, as opposed to directly searching for $\mathbb{I}(\hat{\lambda}_{i},\hat{n})$, noting that $\hat{\lambda}_{i}=1/\lambda_i$. Table~\ref{tbl:searchspace} indicates that the values for $\lambda_i$ can be selected from a range starting at 1.0 (representing straightforward extrapolation) up to $s\times1.25$ (indicating enhanced interpolation compared to PI), using increments of 0.01, with $s$ denoting the intended ratio for expanding the context window.

The parameter $\hat{n}$ determines how many starting token positions keep their original RoPE embedding, foregoing positional interpolation. In practice, we experiment with $\hat{n}$ values chosen from \{0, 1, 2, 4, 8, 12, 16, 20, 24, 28, 32, 64, 128, 256\}. If $\hat{n}=0$, then every token position employs the optimized rescale factors.

\textbf{Evolution-based search}. The search space outlined in Table~\ref{tbl:searchspace} encompasses a wide array of positional interpolation configurations, presenting substantial complexity for systematic exploration. For instance, scaling to $s=4\times$ results in $400^{128/2}\times14$=4$\times10^{167}$ possible options. As the extension ratio increases, the search space grows exponentially. To navigate this vast space, we leverage evolution search~\cite{guo2020single} and incorporate two optimization strategies that significantly enhance search effectiveness. The main steps of the search process are depicted in Algorithm~\ref{alg:search}.

\textit{Efficient initialization of the population}. Rather than relying on random initialization for all $P$ rescale factors, we explicitly include the three RoPE rescale factors associated with PI, NTK, and YaRN as distinct members of the starting population. To complete the set, the other $P-3$ individuals are generated by introducing random mutations to the three rescale factors, each mutation occurring with probability $p$.

\textit{Monotonically non-decreasing constraint}. Once the initial population is produced, we assess each member by calculating its LLM perplexity. This involves applying their respective RoPE rescale factors within the target LLM and measuring perplexity on the input $\mathbf{X}$. The best $k$ individuals are then selected as parents for the next evolutionary step. However, due to the immense size of the search space, conventional mutation and crossover methods may waste resources on suboptimal candidates, resulting in many redundant perplexity evaluations. This inefficiency becomes especially pronounced as $L'$ grows, since each perplexity calculation requires considerable inference time.

To mitigate this, we enforce a monotonic non-decreasing constraint on the sampled RoPE rescaling factors, such that $\lambda_i \leq \lambda_{i+1}$. Only RoPE instances conforming to this constraint are used during perplexity assessment in the LLM, substantially decreasing the computational overhead for searching. In particular, this setup mandates that $\lambda_i$ rises monotonically across the RoPE dimension indices ($i=0,\ldots,63$). The rationale for enforcing dimensional monotonicity stems from NTK theory~\cite{jacot2018neural,tancik2020fourier,ntk}, which indicates that dimensions representing higher frequencies require minimal interpolation (hence, lower values of $\lambda_i$), while those corresponding to lower frequencies benefit from greater interpolation (yielding higher $\lambda_i$ values).

\begin{algorithm}[t]
	\small
	\caption{The search algorithm}
	\textbf{Input:} target LLM, input samples $\mathbf{X}$,   population size $P$, mutation size $N_1$, crossover size $N_2$, max iterations $\mathcal{T}$, mutate probability $p$\\

	\begin{algorithmic}[1]
		\label{alg:search}
		\STATE Top-k $\leftarrow$ $\phi$;	
		\STATE $\text{P}_0$ $\leftarrow$ \textit{Initialize\_population\_with\_optimization} ($P$, $\mathbf{X}$, $p$);
		\FOR{$i = 1$ to $\mathcal{T}$}
		\STATE \textit{Compute\_perplexity} (LLM, $\text{P}_{i-1}$, $\mathbf{X}$);
		\STATE Top-k $\leftarrow$ \textit{Update\_Topk} (Top-k, $\text{P}_{i-1}$);
		\STATE $\text{P}_{mutation}$ $\leftarrow$ \textit{Mutation\_with\_mono\_constraint} (Top-k, $N_1$, $p$);
		\STATE $\text{P}_{crossover}$ $\leftarrow$ \textit{Crossover\_with\_mono\_constraint} (Top-k, $N_2$);
		\STATE $\text{P}_i$ $\leftarrow$ $\text{P}_{mutation}$ $\cup$ $\text{P}_{crossover}$ $\cup$ Top-k;
		\ENDFOR
		\STATE Output the individual in Top-k with minimum perplexity;
	\end{algorithmic}
\end{algorithm}

\textbf{8$\times$ expansion without fine-tuning}. Through evolutionary search, we efficiently determine non-uniform RoPE scaling coefficients, maintaining critical dimensions and positional encoding to reduce information degradation caused by interpolation. Figure~\ref{fig:non-ft} illustrates that our approach successfully increases LLaMA2’s context window from 4k to 32k without requiring fine-tuning. In comparison, other strategies like PI and non-uniform NTK or YaRN result in a sharp rise in perplexity once the window is extended beyond 2$\times$.

\subsection{Extending LLM Context Window to 2048K}
\label{sec:extendto2048k}

\textbf{Progressive extension to 2048k}. In this section, we present our approach for expanding the context window of pre-trained LLMs beyond the conventional 4k limit to over 2048k. Our experiments have shown that employing non-uniform positional interpolation allows for an 8$\times$ expansion of the context length without necessitating additional fine-tuning. However, achieving much greater extensions (such as 512$\times$) requires fine-tuning. A possible approach involves determining optimal RoPE rescaling factors tailored to the 2048k context and subsequently performing fine-tuning. Yet, this procedure is hindered by the immense computational resources demanded for training at such scales. Furthermore, our experience indicates that effectively fine-tuning LLMs for such significant extension ratios poses substantial difficulties (refer to Appendix).

Fortunately, LongRoPE performs well with both pre-trained and fine-tuned variants of the extended LLM. Accordingly, we propose a streamlined, stepwise approach that reaches the desired 2048k sequence length using merely 1k fine-tuning steps, all within a maximum training sequence length of 256k.

$\Diamond$ \textit{LongRoPE search for scaling pre-trained LLMs to 256k context}. Using LLaMA2 as a case study, we perform a search aimed at expanding the context window to 128k and 256k tokens. In this phase, the context is increased by factors of 32$\times$ and 64$\times$, correspondingly.

$\Diamond$ \textit{Fine-tuning to 256k}. Following pre-training, we adapt the large language model to handle a 256k context window through a staged fine-tuning process. Initially, LLaMA2 undergoes 400 fine-tuning steps utilizing RoPE rescaled factors for 128k. Upon completion, the RoPE rescaled factors are swapped for those appropriate to 256k, after which we continue with 600 further fine-tuning steps on the updated checkpoint. This incremental approach shows greater efficiency compared to direct fine-tuning at 256k.

$\Diamond$ \textit{Scaling fine-tuned extended LLM to 2048k using LongRoPE search}. In the last step, we conduct an additional search process on the already fine-tuned LLM with a 256k sequence length. This step successfully achieves an exceptionally large context window of 2048k, all without requiring further rounds of fine-tuning. As a result, the model attains a $512\times$ increase in context window size.

\textbf{Recovery of performance in shorter context windows}.  Upon scaling the context window up to 2048k tokens, we observe degradation in performance specifically within the initial, smaller context window. This drawback is a recognized consequence of positional interpolation~\cite{pi}, where position embeddings corresponding to the original context are compressed into a more confined segment of the embedding space, thereby impairing the model's capabilities. In the scenario where the window is extended by a factor of 512, the positional representations within the native 4k context window are especially densely packed.

To address this issue, we conduct an additional evolutionary search on the expanded LLM to fine-tune RoPE rescale parameters tailored for shorter context ranges, such as 4k and 8k. Since positional interpolation demands are diminished at these lengths, we lower the upper limit for searched $\lambda$. At inference time, the LLM adapts the relevant RoPE rescale values in real-time.

\section{LongRoPE}

\label{sec:method}
Building on these observations, we propose LongRoPE. This approach begins by implementing a streamlined search algorithm that leverages the dual non-uniform patterns, and subsequently applies it to scale the LLM context window to over 2 million tokens.

\subsection{Problem Formulation}

The presence of these two non-uniformities significantly expands the solution space and complicates the optimization process. To mitigate these challenges, we recast the multidimensional non-uniform position interpolation optimization issue as a search problem.

For a LLM targeting a  context window size of $L'$ and lengthy input documents $\mathbf{X}$, where each $\mathbf{x}\in\mathbf{X}$ surpasses $L'$ in token length, we denote the original rotary angle of the $i^{th}$ dimension in RoPE embedding at token position $n$ as  $\frac{n}{\beta^i}$. The optimization problem is then formulated as follows:

\begin{equation}
\begin{gathered}
		\label{eq:problem}

\underset{\mathbf{x} \in \mathbf{X};\, |\mathbf{x}| \geq L'}{\operatorname{arg\,min}}\, \mathcal{L}\left(\text{LLM}\left(\text{RoPE}, \mathbf{X}\right)\right), \text{where}
\\
\scalebox{0.93}{
$\underset {\begin{subarray}{l} i = 0,\dots,\frac{d}{2}-1; \\ n \in [0,|\mathbf{x}|); \end{subarray}}{\text{RoPE}_{(n)}} = {\left[\dots, \cos\left(\mathbb{I}(\hat{\lambda}_i, \hat{n}) \cdot \frac{n}{\beta^i}\right), \sin\left(\mathbb{I}(\hat{\lambda}_i, \hat{n}) \cdot \frac{n}{\beta^i}\right), \dots \right]}$}\\
\text{where } \mathbb{I}(\hat{\lambda}_i, \hat{n}) = \begin{cases}
1 & \text{if } n < \hat{n} \\
\frac{1}{\lambda_i} & \text{if } n \geq \hat{n}
\end{cases}
\end{gathered}
\end{equation}

In this approach, we define a collection of rescale coefficients, $\mathbb{I}(\hat{\lambda}_{i},\hat{n})$, to address two types of non-uniformities. Here, $\hat{\lambda}_i$ represents non-uniformity in RoPE dimensions, while $\hat{n}$ captures non-uniformity in token positions. More precisely, $\mathbb{I}(\hat{\lambda}_{i},\hat{n})$ is employed to modify the rotational angle for the $i^{th}$ dimension of RoPE, using $\hat\lambda_{i}$ as the adjustment coefficient and $\hat{n}$ as the threshold for token positions. For the first $\hat{n}-1$ tokens, $\hat\lambda_{i}$ is not applied, so the conventional RoPE angle $\frac{n}{\beta^i}$ remains unaltered. When $n$ meets or exceeds $\hat{n}$, the rescale factor is incorporated into the computation.

With a specified target context window length $L'$, the goal is to identify the most suitable rescale factors ($\mathbb{I}(\hat{\lambda}_{0},\hat{n})$, $\mathbb{I}(\hat{\lambda}_{1},\hat{n})$, ..., $\mathbb{I}(\hat{\lambda}_{i},\hat{n})$, ...) for each RoPE dimension ranging from $1$ to $d$. By applying these rescale factors to the $\text{RoPE}$ mechanism within the target $\text{LLM}$, we aim to minimize the loss for next-token prediction, denoted by $\mathcal{L}$ (that is, perplexity), when processing input sequences $\mathbf{X}$ of token length $L'$.

\subsection{Searching the Non-uniform Position Interpolation}

To address the challenge posed in Eq.~\ref{eq:problem}, we propose an efficient and straightforward approach that systematically identifies the best RoPE rescale factors. This method is designed to capitalize on the diverse, multidimensional variations present in positional embeddings.

\textbf{Search space}. To accommodate both types of non-uniformity, we construct an expansive search space, outlined in Table~\ref{tbl:searchspace}. In particular, we parameterize a unique rescale factor for each RoPE embedding dimension. For clarity and efficiency in search space formulation, the process involves optimizing over $\lambda_i$ and $\hat{n}$, as opposed to directly targeting $\mathbb{I}(\hat{\lambda}_{i},\hat{n})$, with $\hat{\lambda}_{i}$ defined as $1/\lambda_i$. According to Table~\ref{tbl:searchspace}, $\lambda_i$ spans from a minimum value of 1.0 (corresponding to straightforward extrapolation) up to $s\times1.25$ (which allows more aggressive interpolation than PI), incrementing in steps of 0.01; here, $s$ denotes the ratio by which the context window is extended.

The parameter $\hat{n}$ determines how many of the initial token positions remain unaltered by position interpolation, meaning they utilize the standard RoPE embedding directly. In our experiments, $\hat{n}$ is selected from the set \{0, 1, 2, 4, 8, 12, 16, 20, 24, 28, 32, 64, 128, 256\}. Notably, setting $\hat{n}=0$ implies that every token position employs the learned rescale factors.

\textbf{Evolution-based search}. The search space outlined in Table~\ref{tbl:searchspace} encompasses a wide array of positional interpolation strategies, making efficient navigation particularly demanding. For instance, an extension with $s=4\times$ results in $400^{128/2}\times14$ or 4$\times10^{167}$ potential configurations. As the extension ratio increases, the search space expands at an exponential rate. To cope with these complexities, we adopt evolution search~\cite{guo2020single} and propose two optimizations to substantially enhance the efficiency of the search process. The overall search methodology is depicted in Algorithm~\ref{alg:search}.

\textit{Enhanced initial population setup}. In place of randomly generating $P$ rescale factors for the population, the three RoPE rescale factors associated with PI, NTK, and YaRN are directly incorporated as distinct members at the outset. To create the other $P$-3 population members, each of the three rescale factors is randomly altered with a likelihood of $p$.

\textit{Monotonically non-decreasing constraint}. Once the initial population has been established, we evaluate the LLM perplexity for every individual. This involves applying the relevant RoPE rescale factors to the designated LLM and measuring the perplexity for the input $\mathbf{X}$. The individuals with the best k perplexity scores are chosen as parents for the subsequent evolutionary process. Due to the large size of the search space, straightforward application of mutation and crossover may lead to suboptimal candidates, which in turn results in many redundant perplexity evaluations. This issue becomes more pronounced as $L'$ increases, since each perplexity assessment requires a resource-intensive inference.

To mitigate this issue, we enforce a monotonicity constraint whereby the sampled RoPE rescaling factors must not decrease, i.e., $\lambda_i \leq \lambda_{i+1}$. Only RoPE configurations that adhere to this constraint are utilized when evaluating LLM perplexity, thereby streamlining the search process considerably. Specifically, the requirement stipulates that $\lambda_i$ exhibits a strictly monotonic increase as a function of the RoPE dimension (with $i$ ranging from 0 to 63). This approach is motivated by NTK theory~\cite{jacot2018neural,tancik2020fourier,ntk}, which posits that dimensions with higher frequencies (i.e., lower indices) necessitate minimal interpolation (meaning a smaller $\lambda_i$), while dimensions associated with lower frequencies (i.e., higher indices) are capable of greater interpolation (indicating a larger $\lambda_i$).

\begin{algorithm}[t]
	\small
	\caption{The search algorithm}
	\textbf{Input:} target LLM, input samples $\mathbf{X}$, population size $P$, mutation size $N_1$, crossover size $N_2$, maximum iterations $\mathcal{T}$, mutation probability $p$\\

	\begin{algorithmic}[1]
		\label{alg:search}
		\STATE Top-k $\leftarrow \phi$;	
		\STATE $\text{P}_0 \leftarrow$ \textit{Initialize\_population\_with\_optimization} ($P$, $\mathbf{X}$, $p$);
		\FOR{$i$ = 1 to $\mathcal{T}$}
		\STATE \textit{Compute\_perplexity} (LLM, $\text{P}_{i-1}$, $\mathbf{X}$);
		\STATE Top-k $\leftarrow$ \textit{Update\_Topk} (Top-k, $\text{P}_{i-1}$);
		\STATE $\text{P}_{mutation} \leftarrow$ \textit{Mutation\_with\_mono\_constraint} (Top-k, $N_1$, $p$);
		\STATE $\text{P}_{crossover} \leftarrow$ \textit{Crossover\_with\_mono\_constraint} (Top-k, $N_2$);
		\STATE $\text{P}_i \leftarrow \text{P}_{mutation} \cup \text{P}_{crossover} \cup$ Top-k;
		\ENDFOR
		\STATE Return the member in Top-k yielding the lowest perplexity;
	\end{algorithmic}
\end{algorithm}

\textbf{8$\times$ extension without fine-tuning}. Utilizing an evolutionary search, we strategically select non-uniform RoPE rescale factors that retain crucial dimensions and positions, thereby reducing the information loss that typically occurs during interpolation. As shown in Fig.\ref{fig:non-ft}, this approach successfully expands LLaMA2's context window from 4k to 32k without requiring any fine-tuning. Conversely, alternative techniques like PI and non-uniform NTK and YaRN result in a significant increase in perplexity once the context window is doubled.

\subsection{Extending LLM Context Window to 2048K}
\label{sec:extendto2048k}

\textbf{Progressive extension to 2048k}. In this section, we present our approach for expanding the context window of pre-trained LLMs far beyond the conventional 4k limit, reaching up to 2048k. As shown in our results, leveraging non-uniform positional interpolation enables an 8$\times$ context window extension without requiring additional fine-tuning. However, for far greater expansions—such as 512$\times$—fine-tuning becomes essential. One strategy involves determining RoPE rescaled factors specifically for the target 2048k window, followed by model fine-tuning. Nevertheless, this route is hindered by the substantial computational demands. Furthermore, our observations indicate that fine-tuning LLMs with such a substantial context window increase is notably difficult to optimize effectively (refer to Appendix for details).

Fortunately, LongRoPE demonstrates efficacy with both base models and those subjected to fine-tuning for extended contexts. To this end, we propose a progressive and efficient approach that attains the desired 2048k context window, requiring only 1k fine-tuning steps within a maximum training sequence length of 256k.

$\Diamond$ \textit{Exploring LongRoPE search to scale pre-trained LLMs up to 256k}. Using LLaMA2 as a case study, we perform search experiments with desired context window lengths set at 128k and 256k. At this phase, the model undergoes expansions of 32$\times$ and 64$\times$ its original capacity, respectively.

$\Diamond$ \textit{Fine-tuning to 256k}. The pre-trained LLM is subsequently fine-tuned in order to support a context window of 256k. More concretely, we begin by fine-tuning LLaMA2 for 400 steps utilizing RoPE rescaled factors corresponding to 128k. After this phase, we switch the RoPE rescaled factors to those for 256k on the resulting checkpoint and carry out 600 further fine-tuning steps. Compared to fine-tuning directly to 256k, this staged approach demonstrates greater efficiency.

$\Diamond$ \textit{Expanding fine-tuned extended LLM to 2048k using LongRoPE search}. In the final stage, we carry out a secondary search on the LLM that has already been fine-tuned for 256k sequence length. As a result, the context window expands to an unprecedented 2048k, achieved without any additional fine-tuning steps. This process yields a total extension factor of $512\times$.

\textbf{Restoring performance in shorter context windows}. Upon enlarging the context window to an extensive 2048k tokens, we observe diminished model performance within the initial context length. This phenomenon aligns with previously documented limitations of positional interpolation~\cite{pi}, where increasing the overall window compresses positional embeddings of the original context into a smaller portion of the embedding space, consequently impairing the language model's effectiveness. Notably, with a 512$\times$ increase, tokens within the baseline 4k context window are subject to a significant compression of their corresponding positions.

To address this, we conduct an additional evolution search on the expanded LLM to optimize the RoPE rescale factors specifically for shorter context lengths (such as 4k and 8k). For these shorter lengths, we lower the upper bound on the searched $\lambda$ since less positional interpolation is needed. During inference, the LLM adaptively modifies the relevant RoPE rescale factors.

 \section{Experiments}

\label{sec:exp}
\subsection{Setup}
\textbf{Evaluation Tasks and models}. LongRoPE is implemented on LLaMA2-7B and Mistral-7B, and its effectiveness is assessed in three key areas: (1) measuring perplexity scores of large language models with extended contexts when processing lengthy documents; (2) evaluating performance on the Passkey retrieval task, which tests the capability to extract a specified passkey from among abundant extraneous content; and (3) analyzing results on standard LLM benchmark datasets confined to a 4096-token context window.

\textbf{Fine-tuning}. For LLaMA2, the training protocol involves a learning rate of 2e-5, implemented with a linear decay schedule, and a global batch size set at 32. We conduct fine-tuning over 400 iterations on the Redpajama~\cite{redpajama} dataset, which is partitioned into 128k token segments, each marked by BOS and EOS tokens at the boundaries. After completing these steps, the resulting checkpoint is used as a starting point for an additional 600 training steps to extend the context window to 256k. Training at the 128k context length utilizes 8 A100 GPUs with the distributed system~\cite{cube}, whereas scaling up to 256k context demands 16 A100 GPUs. For Mistral, the learning rate remains constant at 1e-6 with a global batch size of 64. Both the 128k and 256k models are trained following the YaRN~\cite{yarn} configuration: 400 steps on Together Computer's Long-Data Collections~\cite{mistraldata}, each with a sequence length of 16k. Four A100 GPUs are allocated for this training process.

\textbf{Search}. When the target window size is at most 256k, the following configuration is adopted: $P$=64, $N_1$=$N_2$=16, $p$=0.3, $\mathcal{T}$=40, and the top-32 candidates are chosen for mutation and crossover per iteration. Perplexity evaluation is conducted on 5 randomly drawn samples from the PG19 validation set, each meeting the specified minimum target context length. For windows greater than 512k, population size, as well as the numbers for mutation and crossover, are reduced by half. In this regime, perplexity is assessed using 3 random validation samples from Pile-Books3~\cite{pile}.

\textbf{Baselines}. To achieve a 2048k context window, models were fine-tuned on both 128k and 256k sequence lengths. This process results in LongRoPE-2048k (ft=128k) for LLaMA2 and LongRoPE-2048k (ft=256k) for Mistral. The four resulting models are evaluated against leading context window extension baselines, specifically leveraging open-source LLMs that have been fine-tuned subsequent to applying positional interpolation methods, such as PI, NTK, and YaRN. The baseline models considered include Together-32k~\cite{together}, Code LLaMA~\cite{codellama}, LongLoRA-full-FT-100k~\cite{longlora}, YaRN-LLaMA, and YaRN-Mistral~\cite{yarn}.

\subsection{Main Results}

\label{sec:mainresult}
\textbf{Language modeling over extended sequences up to 256k}. Our analysis first centers on benchmarking against leading extended LLMs, focusing on evaluation sequences of 256k tokens. To illustrate the breadth of our approach, we report results across two datasets: Proof-pile~\cite{pg19} and the PG19~\cite{pile} test sets. Perplexity is computed over a range of context lengths via a sliding window protocol of 256 tokens. For PG19, the entirety of its 100 test documents is used. Proof-pile evaluations adhere to the YaRN~\cite{yarn} methodology, involving random selection of 10 samples, each exceeding 128k tokens in length.

Table~\ref{tbl:proof-pile} and Table~\ref{tbl:pg19} present a comparison of perplexity scores for LLaMA2 and Mistral models extended with various interpolation techniques on the Proof-pile and PG19 datasets, respectively. Two primary findings emerge: \textbf{(1)} the extended models consistently exhibit lower perplexity as the evaluation length increases from 4k up to 256k, indicating effective utilization of longer context. \textbf{(2)} Notably, despite having context windows expanded to 16$\times$ their original length—a scenario in which models often struggle to retain short-context performance—our LongRoPE-2048k variants surpass current leading baselines at the 256k context length.

\textbf{Long sequence language modeling for inputs over 2000k}. To assess performance on exceptionally lengthy texts, the Books3~\cite{pile} dataset is employed. For efficiency during evaluation, we randomly choose 20 books that each contain more than 2048k tokens, applying a sliding window of 256k tokens for analysis.

As indicated in Table~\ref{tbl:books3}, LongRoPE enables both LLaMA2-7B and Mistral-7B to attain a context window size of 2048k, while maintaining perplexity metrics that are either on par with or superior to competing approaches across shorter windows ranging from 8k to 128k. The results highlight a significant distinction between LLaMA2 and Mistral when evaluated at 2048k. Mistral demonstrates stronger performance over shorter sequence lengths but its perplexity surpasses 7 after 256k. LLaMA2, as expected, exhibits decreasing perplexity as the context window enlarges, with only slight rises at the 1024k and 2048k marks. Additionally, LongRoPE-2048k on LLaMA2 achieves superior outcomes when fine-tuned at 256k compared to 128k, which can be attributed to its lower secondary extension ratio (8$\times$ instead of 16$\times$). By contrast, Mistral obtains more favorable results when fine-tuning is conducted at a window length of 128k. This observation stems primarily from the adoption of YaRN's protocol, wherein the training length remains at 16k for both 128k and 256k fine-tuning scenarios, thereby influencing Mistral's capacity for subsequent context window expansion following fine-tuning.

\textbf{Passkey retrieval}. Next, we analyze how context window length impacts generation tasks. Building upon the synthetic evaluation framework for passkey retrieval introduced by ~\cite{passkey}, we assess the model's ability to extract a randomly embedded passkey—a five-digit numeric string—from an extended document. The prompt format utilized in this evaluation is described in the appendix. For this assessment, the passkey is inserted at a random position throughout the evaluation context, selected uniformly, and the retrieval task is repeated for 10 different trials.

Fig.~\ref{fig:passkey} illustrates the comparison of retrieval accuracy against various baseline approaches. The performance of previously established models declines precipitously, reaching zero accuracy when sequence length exceeds 128k. In stark contrast, LongRoPE-LLaMA2-2048k (ft=256k) sustains a robust retrieval accuracy ($\ge$90\%) across an extensive range from 4k to 2048k tokens, even in the context of the demanding task of extracting a passkey from sequences containing millions of tokens. Additionally, LongRoPE-Mistral-2048k (ft=128k) consistently achieves 100\% accuracy until 1800k tokens, after which accuracy decreases to 60\% at 2048k. This trend is consistent with the findings presented in Table~\ref{tbl:books3}, where a modest rise in perplexity is observed at the 2048k token length.

\textbf{Evaluation on established benchmarks in the original context window}. We assess the performance of LongRoPE-2048k models on their native context length utilizing the Hugging Face Open LLM Leaderboard~\cite{huggingfaceboard}, employing both zero-shot and few-shot approaches. Specifically, the benchmarks comprise ARC-Challenge~\cite{allenai:arc} with 25-shot setting, HellaSwag~\cite{zellers2019hellaswag} in 10-shot mode, MMLU~\cite{mmlu} with 5-shot configuration, and TruthfulQA~\cite{lin2021truthfulqa} under zero-shot evaluation.

Table~\ref{tbl:commonsense} demonstrates that our models perform similarly to others on the initial benchmark, which was intended for use with shorter context windows. Notably, our approach exceeds the baseline Mistral on TruthfulQA by a margin of +0.5\%. While LongRoPE-LLaMA2-2048k, fine-tuned with a 256k context length, experiences marginally greater declines in performance, its results are still acceptable for the majority of tasks.

\subsection{Ablation Results}

\textbf{Assessing the efficacy of the second positional interpolation}. Within our progressive extension framework, our search algorithm enables a secondary, non-uniform positional interpolation applied to the fine-tuned, length-extended LLMs. To examine the utility of this approach, we implement it on our fine-tuned LLaMA2-256k model. We subsequently extend this model to sequence lengths of 512k, 1024k, and 2048k using both PI and YaRN techniques. As illustrated in Table~\ref{tbl:secondsearch}, our method of non-uniform positional interpolation maintains perplexity at a stable level. Conversely, when using PI and YaRN, perplexity escalates sharply as the extension ratio grows.

\textbf{Recovery performance at reduced context lengths.} To address the degradation in effectiveness observed at shorter contexts, we refine the RoPE factors for LongRoPE-2048k through our search process. In particular, we constrain the maximum permissible scale factors during search to favor reduced interpolation when handling brief 4k and 8k sequence lengths. Table~\ref{tbl:shortperformance} reports the perplexity measures for LongRoPE-LLaMA2-2048k evaluated on Proof-pile at these context sizes, as well as mean LLM benchmark accuracy. The findings indicate notable gains in performance for models operating at shorter contexts.

\textbf{Assessment of the two non-uniformity types}. To determine the individual contributions of each non-uniformity aspect, we conduct an ablation study. This involves two distinct experiments: (i) expanding LLaMA2-7B to shorter sequence lengths (16k and 32k) using three strategies—PI, dimension-only RoPE search, and a combined search for both non-uniformities; (ii) increasing the length of our 256k fine-tuned LLaMA2 model to 2048k using analogous methods. Perplexity scores are obtained without any additional fine-tuning. According to Table~\ref{tbl:searchdimension}, adjusting the RoPE dimension non-uniformity yields a notable decrease in perplexity compared to the linear interpolation applied in PI. Token position non-uniformity notably enhances results at 16k and 32k, though its influence wanes at the 2048k length, possibly owing to the extreme sequence length. Simply retaining original tokens without any interpolation offers minimal benefit in this setting, which we identify as a topic for further investigation.

\section{Related Works}

Alongside strategies relying on position interpolation, this section reviews alternative methods proposed in related literature.

\textbf{Retrieval-based} approaches employ external memory mechanisms to store extended historical context, accompanied by retrieval components to access pertinent documents during inference~\cite{longllama,wang2023augmenting,borgeaud2022improving}. Typically, these frameworks require significant alterations to the architecture of large language models (LLMs). Conversely, our methodology is comparatively streamlined, involving only minor adjustments to positional embeddings. Furthermore, our approach demonstrates versatility across a range of long context tasks beyond retrieval scenarios, including comprehensive document summarization and few-shot learning.

\textbf{Attention-based context window extensions}. In addition to methods that interpolate positional embeddings, other studies have leveraged modifications to attention mechanisms to enable longer input context windows within the same native LLM context size~\cite{han2023lminfinite, streamingllm,parallelwindow}. Central to these approaches is the use of innovative attention masks that address the problem of increasing attention overhead from additional positions. Notably, such strategies can be combined effectively with positional interpolation techniques.

\textbf{Fine-tuning based approaches} investigate strategies for adapting pre-trained LLMs with adjusted position embeddings to accommodate longer sequences. Studies such as Code LLaMA~\cite{codellama}, LLaMA2 Long~\cite{xiong2023effective}, and ScaledRoPE~\cite{liu2023scaling} employ an enlarged base value for RoPE and then perform fine-tuning on the desired context length. Our approach provides adaptability across different target sequence lengths and extends capabilities to beyond 2M tokens. With recent advancements, given that fine-tuning for extended contexts (exceeding 128k tokens) requires heavy GPU investment, methods such as LongLoRA~\cite{longlora} and PoSE~\cite{zhu2023pose} have emerged to reduce computational costs. Our proposed technique operates independently and can be integrated with these efficient fine-tuning solutions.

\section{Conclusion}

In this study, we introduce LongRoPE, a technique that significantly increases the context length available to LLMs up to a groundbreaking 2048k, all while preserving their effectiveness within the original, shorter context window. Our approach leverages two distinct non-uniformities found in RoPE positional embeddings through a computationally efficient evolutionary search process. This methodology yields dual advantages: it supplies an advantageous starting point for subsequent fine-tuning and allows for an 8$\times$ expansion of the context window in the absence of fine-tuning. Additionally, we present a progressive extension method utilizing LLMs fine-tuned on 256k-length contexts to achieve a 2048k context window without incurring additional fine-tuning steps. Thorough experimentation demonstrates the efficacy of LongRoPE. We anticipate that LongRoPE-2048k models will unlock a wide array of long-context use cases and stimulate further exploration in the field.

\section*{Broader Impacts} The research detailed in this paper aims to contribute to the progression of Machine Learning. While our work may lead to various societal outcomes, we do not identify any particular consequences that warrant special attention at this time.

	\balance

\end{document}